% 原书 p.3–p.6（PDF p.11–p.14）
% >>> p011
\section{辛流形的例子（Examples of Symplectic Manifolds）}

\begin{keytakeaways}
\begin{itemize}
  \item 四个基本模型：欧氏空间 $(\R^{2n},\omega_{\mathrm{std}})$、余切丛 $(T^*Q,\omega_{\mathrm{can}})$、
        闭曲面 $(\Sigma_g,\omega_{\mathrm{area}})$、复射影空间 $(\CP^n,\omega_{\mathrm{FS}})$；
        积运算 (1.4) 则给出无穷多新例子。
  \item 数学工具：典范原形式 $\lambda_{\mathrm{can}}$ 与无坐标表达式 (1.3)、$\omega^n=\vol$、
        $\dd\bar\partial$-引理式的构造 (1.5)。
  \item 本节结论：Darboux 定理与 Weinstein 定理说明辛几何没有局部不变量（局部只剩标准形），
        于是这门学科的重心完全落在整体（拓扑）性质上。
  \item 定量差别：$S^2\times S^2$ 上的 $\omega_{\mathrm{prod}}$ 与 $\omega_{\mathrm{deformed}}$ 在辛意义上
        本质不同（Gromov），这是「辛形式不只是体积形式」的第一个明确信号。
\end{itemize}
\end{keytakeaways}

本节列举辛流形的一些重要例子。

\begin{wexam}{1.2.1}（欧氏空间）
取 $\R^{2n}$，带坐标 $(x_1,y_1,\cdots,x_n,y_n)$ 及如下闭二形式
\begin{equation}
  \omega_{\mathrm{std}}:=\sum_{i=1}^{n}\dd x_i\wedge\dd y_i .
  \tag{1.1}
\end{equation}
计算可得 $\omega_{\mathrm{std}}^n=C(n)\cdot\dd x_1\wedge\dd y_1\wedge\cdots\wedge\dd x_n\wedge\dd y_n$，
其中 $C(n)$ 为某个正常数。因此，$(\R^{2n},\omega_{\mathrm{std}})$ 是辛流形。此外，对任意开子集
$U\subseteq\R^{2n}$，$(U,\omega_{\mathrm{std}}|_U)$ 也是辛流形。\hfill$//$
\end{wexam}

辛流形 $(\R^{2n},\omega_{\mathrm{std}})$ 是\textbf{恰当辛流形}（exact symplectic manifold）的标准例子，
这里「恰当」是指其辛结构 $\omega_{\mathrm{std}}$ 是一个恰当二形式。例如 $\omega_{\mathrm{std}}=\dd\lambda_{\mathrm{std}}$，
其中
\[
  \lambda_{\mathrm{std}}=\frac{1}{2}\sum_{i=1}^{n}(x_i\,\dd y_i-y_i\,\dd x_i).
\]
这个例子的重要性体现在下面这个基本定理中，其严格陈述与证明将在第 4 章 4.3 节给出。
% [Tutor Note: 原文此处写作「Section 4.3 in Chapter 5」（第 5 章 4.3 节），但按本书记号「4.3 节」指
% 第 4 章第 3 节；查原书目录，第 4 章为 Moser's Argument，其 4.3 节恰为 Neighborhood Theorems
% （原书印刷页 92），而第 5 章为 Floer Theory，其 5.3 节是 Homological Rigidity。故「第 5 章」系
% 笔误，译文按目录证据更正为第 4 章，并记入 review_ch01.md。]

% >>> p012
\begin{wthm}{1.2.1}（Darboux）
任何辛流形 $(M^{2n},\omega)$ 都局部地（在辛的意义下）等价于 $(U,\omega_{\mathrm{std}}|_U)$，其中
$U\subset\R^{2n}$ 为某个开子集。
\end{wthm}

由于定理 1.2.1，辛流形没有有趣的局部不变量，因此辛几何研究的是辛流形的整体性质，这也为该学科
引出了一个更准确的名称：辛拓扑（symplectic topology）。

\begin{wexam}{1.2.2}（余切丛）
设 $Q$ 为微分 $n$ 维流形 $Q$。其余切丛 $T^*Q$ 具有典范辛结构，记作 $\omega_{\mathrm{can}}$。
描述这一辛结构最方便的方式是使用局部坐标卡。具体地，设 $\pi:T^*Q\to Q$ 由 $(q,p)\mapsto q$ 定义，
为自然投影映射。取 $Q$ 的任一开坐标卡 $U$，使得原像 $\pi^{-1}(U)$ 具有局部坐标
$(q_1,\cdots,q_n,p_1,\cdots,p_n)$，其中 $p_i$ 表示纤维中的余向量。于是在 $\pi^{-1}(U)$ 上，
辛结构 $\omega_{\mathrm{can}}$ 写作
\begin{equation}
  \omega_{\mathrm{can}}=\sum_{i=1}^{n}\dd q_i\wedge\dd p_i .
  \tag{1.2}
\end{equation}
与 (1.1) 中欧氏空间上的 $\omega_{\mathrm{std}}$ 类似，(1.2) 中 $T^*Q$ 上的 $\omega_{\mathrm{can}}$
实际上也是恰当的，即 $\omega_{\mathrm{can}}=\dd\lambda_{\mathrm{can}}$。此外，典范原形式
$\lambda_{\mathrm{can}}$ 具有一个无坐标的表达式：对任意向量 $Y\in T_{(q,p)}T^*Q$，原形式
$\lambda_{\mathrm{can}}$ 由下式定义
\begin{equation}
  (\lambda_{\mathrm{can}})_{(q,p)}(Y):=p\big((\pi_{(q,p)})_*(Y)\big).
  \tag{1.3}
\end{equation}
容易验证，在上述局部坐标下，$\lambda_{\mathrm{can}}|_{\pi^{-1}(U)}=\sum_{i=1}^{n}p_i\,\dd q_i$。\hfill$//$
\end{wexam}
% [Tutor Note: 原文疑误——「设 Q 为微分 n 维流形 Q」句末重复了符号 Q（应为「设 Q 为微分 n 维流形」），
% 译文照录并在此说明。记号说明：$\pi_{(q,p)}$ 表示投影 $\pi$ 在点 $(q,p)$ 处取值的映射，
% $(\pi_{(q,p)})_*$ 是它在切空间上的微分；$p$ 在此处是 $T_q^*Q$ 中的余向量，故 $p(\cdot)$ 是余向量对
% 切向量的自然配对，与「点 $p$」的用法（如 $p\in M$）不同，阅读时须分清。另注意 (1.2) 的符号顺序
% $\dd q_i\wedge\dd p_i$：若改用 $\dd p_i\wedge\dd q_i$，则典范原形式变为 $-\sum p_i\dd q_i$，
% 这是文献中常见的差别来源。]

\begin{wexer}{1.2.1}
给定光滑流形 $Q$ 与微分同胚 $f:Q\to Q$，证明 $f$ 提升为微分同胚 $\varphi_f:T^*Q\to T^*Q$，满足
$(\varphi_f)^*\lambda_{\mathrm{can}}=\lambda_{\mathrm{can}}$，其中 $\lambda_{\mathrm{can}}$ 是余切丛
$T^*Q$ 上典范辛结构 $\omega_{\mathrm{can}}=\dd\lambda_{\mathrm{can}}$ 的典范原形式。
\end{wexer}

再次与欧氏空间的情形类似，余切丛在辛几何中的重要性来自下面这个基本定理（其严格陈述与证明将在
4.3 节给出）。

\begin{wthm}{1.2.2}（Weinstein）
对某个子流形 $L\subset(M,\omega)$（称为\textbf{拉格朗日子流形}，将在 1.6 节介绍），它在 $M$ 中
具有一个邻域，该邻域辛等价于 $(U,\omega_{\mathrm{can}}|_U)$，其中 $U\subset T^*L$ 是余切丛
$T^*L$ 中零截面的一个邻域。
\end{wthm}
% [Tutor Note: 原文此处写作「$U\subset T^*M$ 是余切丛 $T^*L$ 中零截面的邻域」，同一个 $U$ 既在 $T^*M$
% 又在 $T^*L$ 中，前后矛盾；由 $\omega_{\mathrm{can}}$ 只定义在 $T^*L$ 上可知应为 $U\subset T^*L$。
% 译文按更正后的写法排印，详见 review_ch01.md。]

% >>> p013
由于定理 1.2.2，拉格朗日子流形的许多性质——尤其是局部性质——在其余切丛中研究起来更为有效。

\begin{wexam}{1.2.3}（闭曲面）
记 $\Sigma_g$ 为亏格 $g$ 的闭可定向曲面。它是辛流形，其辛结构取为体积形式（面积形式）
$\omega_{\mathrm{area}}$。$\Sigma_g$ 上两个不同的辛结构只不过是彼此的一个重新标度。\hfill$//$
\end{wexam}
% [Tutor Note: 严格地说，「相互只差一个伸缩」应在「相差一个辛同胚」的意义下理解：对同一个 $\Sigma_g$
% 上的两个面积形式 $\omega_1,\omega_2$，一般并不存在常数 $c$ 使 $\omega_1=c\,\omega_2$ 逐点成立；
% 正确说法是：存在微分同胚 $\varphi$ 与常数 $c>0$ 使 $\varphi^*\omega_2=c\,\omega_1$（当总面积相等时
% Moser 稳定性甚至给出 $c=1$ 的辛同胚）。这正说明二维情形辛结构与体积结构没有差别，
% 而差别从四维才开始出现。]

\begin{wexam}{1.2.4}（积）
设 $(M^{2m},\omega_M)$ 与 $(N^{2n},\omega_N)$ 为两个辛流形，$\pi_M,\pi_N$ 分别为 $M\times N$ 到 $M$
与 $N$ 的自然投影映射。可以如下定义 $M\times N$ 上的辛结构
\begin{equation}
  \omega_{\mathrm{prod}}:=\pi_M^*\omega_M+\pi_N^*\omega_N .
  \tag{1.4}
\end{equation}
事实上，$\omega_{\mathrm{prod}}^{m+n}=\binom{m+n}{m}\cdot\left(\pi_M^*\omega_M^m\wedge\pi_N^*\omega_N^n\right)$，
它是 $M\times N$ 上的体积形式。\hfill$//$
\end{wexam}

\begin{wexer}{1.2.2}
设 $(M^{2m},\omega_M)$ 与 $(N^{2n},\omega_N)$ 为两个辛流形。证明：对任意非零的 $s,t$，\textbf{加权积}
（weighted product）
\[
  \omega_{(s,t)}=s\,\pi_M^*\omega_M+t\,\pi_N^*\omega_N
\]
是 $M\times N$ 上的辛结构。
\end{wexer}

作为练习 1.2.2 的推论，积 $S^2\times S^2$ 至少可以被赋予三种辛结构：
\begin{enumerate}[label=(\roman*)]
  \item $(S^2\times S^2,\ \omega_{\mathrm{prod}}=\pi_M^*\omega_{\mathrm{area}}+\pi_N^*\omega_{\mathrm{area}})$；
  \item $(S^2\times S^2,\ \omega_{\mathrm{twisted}}:=\pi_M^*\omega_{\mathrm{area}}-\pi_N^*\omega_{\mathrm{area}})$；
  \item $(S^2\times S^2,\ \omega_{\mathrm{deformed}}:=\pi_M^*\omega_{\mathrm{area}}+(1+\epsilon)\pi_N^*\omega_{\mathrm{area}})$，
        其中 $\epsilon>0$ 充分小。
\end{enumerate}
M. Gromov 的一个深刻结果指出：在 $S^2\times S^2$ 上，$\omega_{\mathrm{deformed}}$ 的某些辛性质与
$\omega_{\mathrm{prod}}$ 有本质的不同（[99]，进一步的研究见 [10]）。

\begin{wexam}{1.2.5}（复射影空间）
回忆 $\CP^n$ 可由以下标准局部坐标卡覆盖
\[
  \begin{aligned}
    U_i&=\big\{[z_0,\cdots,z_n]\ \big|\ (z_0,\cdots,z_n)\in\C^{n+1}\ \text{且}\ z_i\neq0\big\}\\
       &\cong\big\{(w_1,\cdots,w_n)\ \big|\ w_i\in\C\big\}=\C^n .
  \end{aligned}
\]
上面的微分同胚由以下两步复合给出：先把 $[z_0,\cdots,z_n]$ 映到
$\left(\frac{z_0}{z_i},\cdots,1,\cdots,\frac{z_n}{z_i}\right)$，然后去掉 $1$ 并把每一项重新标号为 $w_i$。
在上述等同的意义下，在 $U_i$ 上定义二形式，记作 $\omega_{\mathrm{FS}}|_{U_i}$，如下
\begin{equation}
  \omega_{\mathrm{FS}}|_{U_i}=\frac{\sqrt{-1}}{2}\partial\bar\partial\log(|w|^2+1),
  \qquad\text{其中 } w=(w_1,\cdots,w_n)\in\C^n .
  \tag{1.5}
\end{equation}
然后 $\omega_{\mathrm{FS}}|_{U_i}$ 可用来构造 $\CP^n$ 上的辛结构，其构造步骤概括在下面的练习中。\hfill$//$
\end{wexam}

% >>> p014
\begin{wexer}{1.2.3}
设 $\omega_{\mathrm{FS}}$ 为如上由 (1.5) 所定义的二形式。
\begin{enumerate}[label=(\roman*)]
  \item 证明 $\omega_{\mathrm{FS}}$ 在 $\CP^n$ 上是良定义的。换言之，对任意
        $i\neq j\in\{0,\cdots,n\}$，有
        $(\omega_{\mathrm{FS}}|_{U_i})|_{U_i\cap U_j}=(\omega_{\mathrm{FS}}|_{U_j})|_{U_i\cap U_j}$。
        这个二形式称为 $\CP^n$ 上的 \textbf{Fubini--Study 形式}。
  \item 证明 $\omega_{\mathrm{FS}}$ 是辛结构。
  \item 证明 $\displaystyle\int_{\CP^1}\omega_{\mathrm{FS}}=\pi$（回忆
        $\displaystyle\int_{S^2}\omega_{\mathrm{area}}=4\pi$）。
\end{enumerate}
\end{wexer}
% [Tutor Note: (1.5) 中的 $\partial\bar\partial$ 是复外微分，而 $\sqrt{-1}/2$ 这一常数决定了归一化：
% 在 $\CP^1$ 上取仿射坐标 $z=w_1$，利用 $\Delta\log(1+|z|^2)=4/(1+|z|^2)^2$ 与
% $\partial\bar\partial f=\frac14\Delta f\,\dd z\wedge\dd\bar z$，可算出
% $\int_{\CP^1}\omega_{\mathrm{FS}}=\pi$，恰为 (iii) 所述。文献中更常见的是除以 $2\pi$ 的归一化
% $\frac{\sqrt{-1}}{2\pi}\partial\bar\partial\log(1+|w|^2)$，那时 $\int_{\CP^1}\omega_{\mathrm{FS}}=1$，
% 即 $[\omega_{\mathrm{FS}}]$ 是整上同调类且恰为第一陈类 $c_1$ 的 $1/(2\pi)$ 倍（视约定而定）；
% 本书的这个约定使体积与面积多出因子 $\pi$，与 [99]、[153] 等文献对照公式时须先核对这一常数。]
